\paragraph{Writing two facts in a specified order.}
We gave a pictured novel passage and two questions to a single causal writer, together with banks of plausible excerpts. Both annotated facts appeared among the excerpts on 168 of 384 novel requests. The writer gave an ordered, stopped reply on four of those 168. Across all 384 requests, both the ranked and shifted-excerpt fits gave six such replies. Neither fit passed both question orders for any novel passage in either language. On 12 separately authored pages, ranked excerpts yielded one successful pair of orderings across 24 page--language pairs (Fig.~\ref{fig:ordered-book-replies}a,b). On the first new page the Chinese reply repeated the name of the book repairer and omitted the binding material. The English questions elicited a complete Chinese statement of both facts on that same page (Fig.~\ref{fig:ordered-book-replies}c).

We then trained a clause writer to put one physical insertion mark in its generated wording. The runtime replaced that mark with the reader's selected excerpt and joined the two clauses in question order. This increased successful pairs of orderings from zero to 15 of 192 novel page--language pairs and from one to three of 24 on the authored pages. On a further 64 novel pages the clause writer completed 14 of 128 page--language pairs. Merely listing the two selected excerpts gave 37, nine and 24 successful pairs on the respective sets. Supplying the correct excerpts to the clause writer raised the 96-page result from 15 to 106 pairs. Both evidence selection and surrounding language thus contributed errors. Ordering in this architecture follows the input slots by construction.

A close reading that concealed system identity credited complete supported statements on four of 48 authored requests for each trained writer and zero for the bare excerpt list. The local text Qwen3-8B comparison received 34 credits. The direct writer's four successes included two English requests. All four clause-writer successes used Chinese requests. The clause writer distinguished the tree names used in a printed interview and its first draft in both orders (Fig.~\ref{fig:ordered-book-replies}c). Elsewhere it assigned an instrument to the musician who played it during rehearsal after the passage explicitly described an exchange for the performance. An automated Qwen grader credited 22 direct replies, including a reply that merely repeated the wrong performance roles. The blind close reading was model-based, with prior familiarity with some direct-writer examples. These measurements establish a gain in ordered fact inclusion and individual supported statements, while complete-answer accuracy on the 48 requests remains unchanged.
